% Part: first-order-logic
% Chapter: syntax-and-semantics
% Section: assignments

\documentclass[../../../include/open-logic-section]{subfiles}

\begin{document}

\olfileid[hi]{fol}{syn}{ass}

\olsection{चर-नियतन}

\begin{explain}
कोई !!{variable} नियतन~$s$ \emph{प्रत्येक} चर को मान देता है---और चर अनंत रूप से
अनेक हैं। निस्संदेह यह आवश्यक नहीं है। हम !!{variable} नियतनों से सभी !!{variable}s
को मान देने की अपेक्षा केवल इसलिए करते हैं कि इससे काम बहुत सरल हो जाता है।
किसी पद~$t$ का मान और यह कि कोई !!a{formula}~$!A$,~$s$ के सापेक्ष किसी
!!a{structure} में संतुष्ट है या नहीं, केवल उन मानों पर निर्भर करते हैं जिन्हें
~$s$,~$t$ में आने वाले !!{variable}s और~$!A$ के मुक्त !!{variable}s को नियत करता
है। अगली दो प्रतिज्ञप्तियों का विषय यही है। ``निर्भर करता है'' की धारणा को
सुनिश्चित रूप देने के लिए हम दिखाते हैं कि~$t$ में आने वाले सभी चरों पर सहमत
कोई भी दो चर-नियतन एक ही मान देते हैं; और यदि दो चर-नियतन~$!A$ के सभी मुक्त
चरों पर सहमत हों, तो~$!A$ एक के सापेक्ष तभी और केवल तभी संतुष्ट है जब वह
दूसरे के सापेक्ष संतुष्ट है।
\end{explain}

\begin{prop}\ollabel{prop:valindep}
यदि पद~$t$ के !!{variable}, $x_1$, \dots,~$x_n$ में से हों, और
$s_1(x_i) = s_2(x_i)$ हो, $i = 1$, \dots,~$n$ के लिए, तो
$\Value{t}{M}[s_1]
= \Value{t}{M}[s_2]$।
\end{prop}

\begin{proof}
~$t$ की जटिलता पर आगमन से। आधार प्रकरण में~$t$ कोई !!a{constant}, अथवा
$x_1$, \dots,~$x_n$ में से कोई चर हो सकता है। यदि $t
= c$, तो
$\Value{t}{M}[s_1] = \Assign{c}{M} = \Value{t}{M}[s_2]$। यदि $t = x_i$,
तो प्रतिज्ञप्ति की परिकल्पना से $s_1(x_i) = s_2(x_i)$, और इसलिए
$\Value{t}{M}[s_1] = s_1(x_i) = s_2(x_i) = \Value{t}{M}[s_2]$।

आगमन चरण के लिए मान लें कि $t = \Atom{f}{t_1, \dots, t_k}$ और दावा
$t_1$, \dots, $t_k$ के लिए सत्य है। तब
\begin{align*}
  \Value{t}{M}[s_1] & = \Value{\Atom{f}{t_1, \dots, t_k}}{M}[s_1] \\
  & = \Assign{f}{M}(\Value{t_1}{M}[s_1], \dots, \Value{t_k}{M}[s_1]).
\intertext{$j = 1$, \dots,~$k$ के लिए~$t_j$ के !!{variable},
    $x_1$, \dots,~$x_n$ में से हैं। आगमन-परिकल्पना से
    $\Value{t_j}{M}[s_1] = \Value{t_j}{M}[s_2]$। अतः}
\Value{t}{M}[s_1] & = \Value{\Atom{f}{t_1, \dots, t_k}}{M}[s_1] \\
  & = \Assign{f}{M}(\Value{t_1}{M}[s_1], \dots, \Value{t_k}{M}[s_1]) \\
  & = \Assign{f}{M}(\Value{t_1}{M}[s_2], \dots, \Value{t_k}{M}[s_2]) \\
  & = \Value{\Atom{f}{t_1, \dots, t_k}}{M}[s_2] = \Value{t}{M}[s_2].
\end{align*}
\end{proof}

\begin{prop}\ollabel{prop:satindep}
यदि~$!A$ के मुक्त !!{variable}, $x_1$, \dots,~$x_n$ में से हों, और
$s_1(x_i) = s_2(x_i)$ हो, $i = 1$, \dots,~$n$ के लिए, तो
$\Sat{M}{!A}[s_1]$ तभी और केवल तभी जब $\Sat{M}{!A}[s_2]$।
\end{prop}

\begin{proof}
हम~$!A$ की जटिलता पर आगमन करते हैं। आधार प्रकरण में, जहाँ~$!A$ परमाण्विक
है,~$!A$ निम्न में से हो सकता है:
\iftag{prvTrue}{$\ltrue$,}{}
\iftag{prvFalse}{$\lfalse$,}{}
$\Atom{R}{t_1, \dots, t_k}$, किसी $k$-स्थानी विधेय~$R$ और पदों
$t_1$, \dots,~$t_k$ के लिए, अथवा $\eq[t_1][t_2]$, पदों~$t_1$ और~$t_2$
के लिए। अंतिम दो प्रकरणों में हम !!{biconditional} की केवल अग्र
दिशा दिखाते हैं, क्योंकि विलोम दिशा का प्रमाण सममित है।

\begin{enumerate}
\tagitem{prvTrue}{%
  \indcase{!A}{\ltrue}{$\Sat{M}{\indfrm}[s_1]$ और
    $\Sat{M}{\indfrm}[s_2]$, दोनों।}}{}

\tagitem{prvFalse}{%
  \indcase{!A}{\lfalse}{$\Sat/{M}{!A}[s_1]$ और
    $\Sat/{M}{!A}[s_2]$, दोनों।}}{}

\item
  \indcase{!A}{\Atom{R}{t_1, \ldots, t_k}}{मान लें कि
    $\Sat{M}{\indfrm}[s_1]$। तब
    \[
    \langle \Value{t_1}{M}[s_1], \ldots, \Value{t_k}{M}[s_1] \rangle
    \in \Assign{R}{M}.
    \]
    $i = 1$, \dots,~$k$ के लिए, \olref{prop:valindep} से
    $\Value{t_i}{M}[s_1] =
    \Value{t_i}{M}[s_2]$। अतः हमारे पास
    $\langle \Value{t_i}{M}[s_2], \ldots, \Value{t_k}{M}[s_2] \rangle
    \in \Assign{R}{M}$ भी है, और इसलिए $\Sat{M}{\indfrm}[s_2]$।}
\item
  \indcase{!A}{\eq[t_1][t_2]}{मान लें कि $\Sat{M}{\indfrm}[s_1]$।
    तब $\Value{t_1}{M}[s_1] = \Value{t_2}{M}[s_1]$। अतः
    \begin{align*}
      \Value{t_1}{M}[s_2] & = \Value{t_1}{M}[s_1]
      & \text{(\olref{prop:valindep} से)} \\
      & = \Value{t_2}{M}[s_1]
      & \text{(क्योंकि $\Sat{M}{\eq[t_1][t_2]}[s_1]$)}\\
      &= \Value{t_2}{M}[s_2]
      & \text{(\olref{prop:valindep} से),}
    \end{align*}
    अतः $\Sat{M}{\eq[t_1][t_2]}[s_2]$।}
\end{enumerate}

अब मान लें कि $\Sat{M}{!B}[s_1]$ तभी और केवल तभी जब
$\Sat{M}{!B}[s_2]$, उन सभी !!{formula}s~$!B$ के लिए जो~$!A$ से कम जटिल
हैं। आगमन चरण
~$!A$ के मुख्य संकारक द्वारा निर्धारित प्रकरणों के अनुसार चलता है। प्रत्येक
प्रकरण में हम !!{biconditional} की केवल अग्र दिशा दिखाते हैं; विलोम दिशा का
प्रमाण सममित है। परिमाणकों के प्रकरणों को छोड़कर प्रत्येक प्रकरण में हम उन
उप-!!{formula}s~$!B$ पर आगमन-परिकल्पना लगाते हैं जो~$!A$ के हैं।~$!B$ के मुक्त
चर,~$!A$ के मुक्त चरों में से हैं। अतः यदि $s_1$ और $s_2$,~$!A$ के मुक्त
चरों पर सहमत हैं, तो वे~$!B$ के मुक्त चरों पर भी सहमत हैं, और
आगमन-परिकल्पना~$!B$ पर लागू होती है।

\begin{enumerate}
\tagitem{defNot}{}{%
  \iftag{probNot}{%
    \indcase!{!A}{\lnot !B}{}}{%
    \indcase{!A}{\lnot !B}{यदि $\Sat{M}{\indfrm}[s_1]$, तो
      $\Sat/{M}{!B}[s_1]$; अतः आगमन-परिकल्पना से
      $\Sat/{M}{!B}[s_2]$, और इसलिए $\Sat{M}{\indfrm}[s_2]$।}}}

\tagitem{defAnd}{}{%
  \iftag{probAnd}{%
    \indcase!{!A}{!B \land !C}{}}{%
    \indcase{!A}{!B \land !C}{यदि $\Sat{M}{\indfrm}[s_1]$, तो
      $\Sat{M}{!B}[s_1]$ और $\Sat{M}{!C}[s_1]$; अतः आगमन-परिकल्पना से
      $\Sat{M}{!B}[s_2]$ और $\Sat{M}{!C}[s_2]$। इसलिए
      $\Sat{M}{\indfrm}[s_2]$।}}}

\tagitem{defOr}{}{%
  \iftag{probOr}{%
    \indcase!{!A}{!B \lor !C}{}}{%
    \indcase{!A}{!B \lor !C}{यदि $\Sat{M}{\indfrm}[s_1]$, तो
      $\Sat{M}{!B}[s_1]$ अथवा $\Sat{M}{!C}[s_1]$। आगमन-परिकल्पना से
      $\Sat{M}{!B}[s_2]$ अथवा $\Sat{M}{!C}[s_2]$; अतः
      $\Sat{M}{\indfrm}[s_2]$।}}}

\tagitem{defIf}{}{%
  \iftag{probIf}{%
    \indcase!{!A}{!B \lif !C}{}}{%
    \indcase{!A}{!B \lif !C}{यदि $\Sat{M}{\indfrm}[s_1]$, तो
    $\Sat/{M}{!B}[s_1]$ अथवा $\Sat{M}{!C}[s_1]$। आगमन-परिकल्पना से
    $\Sat/{M}{!B}[s_2]$ अथवा $\Sat{M}{!C}[s_2]$; अतः
    $\Sat{M}{\indfrm}[s_2]$।}}}

\tagitem{defIff}{}{%
  \iftag{probIff}{%
    \indcase!{!A}{!B \liff !C}{}}{%
    \indcase{!A}{!B \liff !C}{यदि $\Sat{M}{\indfrm}[s_1]$, तो या तो
      $\Sat{M}{!B}[s_1]$ और $\Sat{M}{!C}[s_1]$, अथवा
      $\Sat/{M}{!B}[s_1]$ और $\Sat/{M}{!C}[s_1]$। आगमन-परिकल्पना से,
      या तो $\Sat{M}{!B}[s_2]$ और $\Sat{M}{!C}[s_2]$, अथवा
      $\Sat/{M}{!B}[s_2]$ और $\Sat/{M}{!C}[s_2]$। दोनों दशाओं में
      $\Sat{M}{\indfrm}[s_2]$।}}}

\tagitem{defEx}{}{%
  \iftag{probEx}{%
    \indcase!{!A}{\lexists[x][!B]}{}}{%
    \indcase{!A}{\lexists[x][!B]}{यदि $\Sat{M}{\indfrm}[s_1]$, तो कोई
      $m \in \Domain{M}$ ऐसा है कि
      $\Sat{M}{!B}[\Subst{s_1}{m}{x}]$। $s_1' = \Subst{s_1}{m}{x}$ और $s_2' =
      \Subst{s_2}{m}{x}$ रखें।~$!B$ के मुक्त चर $x_1$,
      \dots, $x_n$, और~$x$ में से हैं। $s_1'(x_i) = s_2'(x_i)$, क्योंकि
      $s_1'$ और $s_2'$, $x$-परिवर्ती हैं---क्रमशः $s_1$ और~$s_2$ के---और
      परिकल्पना से $s_1(x_i) = s_2(x_i)$। $s_1'(x) = s_2'(x) = m$; यह
      हमारे द्वारा $s_1'$ और~$s_2'$ को परिभाषित करने के ढंग से मिलता है।
      तब आगमन-परिकल्पना~$!B$ और $s_1'$,
      $s_2'$ पर लागू होती है, अतः $\Sat{M}{!B}[s_2']$। इसलिए, चूँकि
      $s_2' = \Subst{s_2}{m}{x}$, कोई $m \in \Domain{M}$ ऐसा है कि
      $\Sat{M}{!B}[\Subst{s_2}{m}{x}]$, और अतः
      $\Sat{M}{\indfrm}[s_2]$।}}}

\tagitem{defAll}{}{%
  \iftag{probAll}{%
    \indcase!{!A}{\lforall[x][!B]}{}}{%
    \indcase{!A}{\lforall[x][!B]}{यदि $\Sat{M}{\indfrm}[s_1]$, तो
      प्रत्येक $m \in \Domain{M}$ के लिए
      $\Sat{M}{!B}[\Subst{s_1}{m}{x}]$। हमें दिखाना है कि प्रत्येक
      $m \in \Domain{M}$ के लिए $\Sat{M}{!B}[\Subst{s_2}{m}{x}]$ भी।
      अतः कोई स्वेच्छ $m \in \Domain{M}$ लें, और
      $s_1' = \Subst{s}{m}{x}$ तथा $s_2' =
      \Subst{s}{m}{x}$ पर विचार
      करें। हमारे पास $\Sat{M}{!B}[s_1']$ है।~$!B$ के मुक्त चर $x_1$,
      \dots, $x_n$, और~$x$ में से हैं। $s_1'(x_i) = s_2'(x_i)$, क्योंकि
      $s_1'$ और $s_2'$, $x$-परिवर्ती हैं---क्रमशः $s_1$ और~$s_2$ के---और
      परिकल्पना से $s_1(x_i) = s_2(x_i)$। $s_1'(x) = s_2'(x) = m$; यह
      हमारे द्वारा $s_1'$ और~$s_2'$ को परिभाषित करने के ढंग से मिलता है।
      तब आगमन-परिकल्पना~$!B$ और~$s_1'$,
      $s_2'$ पर लागू होती है, और हमें $\Sat{M}{!B}[s_2']$ मिलता है। यह
      प्रत्येक $m \in \Domain{M}$ पर लागू है, अर्थात्
      $\Sat{M}{!B}[\Subst{s_2}{m}{x}]$, सभी~$m \in
      \Domain{M}$ के लिए; अतः
      $\Sat{M}{\indfrm}[s_2]$।}}}
\end{enumerate}
आगमन से हमें मिलता है कि $\Sat{M}{!A}[s_1]$ तभी और केवल तभी जब
$\Sat{M}{!A}[s_2]$, जब भी~$!A$ के मुक्त !!{variable} $x_1$, \dots,
$x_n$ में से हों और $s_1(x_i)=s_2(x_i)$ हो, $i = 1$, \dots,~$n$ के लिए।
\end{proof}

\begin{probtag}{probNot,probOr,probAnd,probIf,probIff,probEx,probAll}
\olref[fol][syn][ass]{prop:satindep} का प्रमाण पूरा कीजिए।
\end{probtag}

\begin{explain}
!!^{sentence}s में कोई मुक्त चर नहीं होता; अतः कोई भी दो चर-नियतन किसी
वाक्य के सभी (शून्य) मुक्त चरों को समान वस्तुएँ नियत करते हैं। अभी सिद्ध
प्रतिज्ञप्ति का अर्थ है कि कोई !!a{sentence}, किसी चर-नियतन के सापेक्ष किसी
संरचना में संतुष्ट है या नहीं, यह नियतन से पूर्णतः स्वतंत्र है। हम इस तथ्य
को दर्ज कर लेते हैं। इससे आगे आने वाली, चर-नियतन का उल्लेख किए बिना किसी
!!a{structure} में !!a{sentence} की संतुष्टि की परिभाषा उचित ठहरती है।
\end{explain}

\begin{cor}
\ollabel{cor:sat-sentence}
यदि~$!A$ कोई !!a{sentence} और~$s$ कोई चर-नियतन है, तो
$\Sat{M}{!A}[s]$ तभी और केवल तभी जब $\Sat{M}{!A}[s']$, प्रत्येक चर-नियतन~$s'$
के लिए।
\end{cor}

\begin{proof}
  $s'$ को कोई भी चर-नियतन लें। चूँकि~$!A$ कोई !!a{sentence} है, उसमें
  कोई मुक्त चर नहीं है; अतः प्रत्येक चर-नियतन~$s'$,~$!A$ के सभी मुक्त
  चरों को तुच्छतः वही वस्तुएँ नियत करता है जो~$s$ करता है। इसलिए
  \olref{prop:satindep} की शर्त पूरी है, और हमारे पास
  $\Sat{M}{!A}[s]$ तभी और केवल तभी जब $\Sat{M}{!A}[s']$।
\end{proof}

\begin{defn}
\ollabel{defn:satisfaction}
यदि~$!A$ कोई !!a{sentence} है, तो हम कहते हैं कि !!a{structure}~$\Struct M$,
~$!A$ को \emph{संतुष्ट करती है}, $\Sat{M}{!A}$, तभी और केवल तभी जब
$\Sat{M}{!A}[s]$, सभी चर-नियतनों~$s$ के लिए।
\end{defn}

यदि $\Sat{M}{!A}$, तो हम सरलता से यह भी कहते हैं कि
\emph{$!A$,~$\Struct{M}$ में सत्य है।} संतुष्टि की धारणा अलग-अलग
!!{sentence}s से स्वाभाविक रूप से !!{sentence}s के समुच्चयों तक विस्तृत
होती है।

\begin{defn}
\ollabel{defn:sat}
यदि~$\Gamma$, !!{sentence}s का समुच्चय है, तो इस~$\Gamma$ के लिए हम कहते हैं कि
!!a{structure}~$\Struct M$,~$\Gamma$ को \emph{संतुष्ट करती है},
$\Sat{M}{\Gamma}$, तभी और केवल तभी जब $\Sat{M}{!A}$, सभी $!A \in \Gamma$
के लिए।
\end{defn}

\begin{prop}\ollabel{prop:sentence-sat-true}
  मान लें कि~$\Struct{M}$ कोई !!a{structure},~$!A$ कोई !!a{sentence}, और
  $s$ कोई चर-नियतन है। $\Sat{M}{!A}$ तभी और केवल तभी जब
  $\Sat{M}{!A}[s]$।
\end{prop}

\begin{proof}
अभ्यास।
\end{proof}

\begin{prob}
\olref[fol][syn][ass]{prop:sentence-sat-true} सिद्ध कीजिए।
\end{prob}

\begin{prop}\ollabel{prop:sat-quant}
  मान लें कि~$!A(x)$ में केवल~$x$ मुक्त आता है, और~$\Struct M$ कोई
  !!a{structure} है। तब:
  \begin{tagenumerate}{prvEx,prvAll}
    \tagitem{prvEx}{$\Sat{M}{\lexists[x][!A(x)]}$ तभी और केवल तभी जब
      $\Sat{M}{!A(x)}[s]$, कम-से-कम एक चर-नियतन~$s$ के लिए।}{}
    \tagitem{prvAll}{$\Sat{M}{\lforall[x][!A(x)]}$ तभी और केवल तभी जब
      $\Sat{M}{!A(x)}[s]$, सभी चर-नियतनों~$s$ के लिए।}{}
\end{tagenumerate}
\end{prop}

\begin{proof}
अभ्यास।
\end{proof}

\begin{prob}
\olref[fol][syn][ass]{prop:sat-quant} सिद्ध कीजिए।
\end{prob}

\begin{prob}
\DeclareRobustCommand{\VDash}{\mathrel{||}\joinrel\Relbar}
मान लें कि~$\Lang L$ कोई ऐसी भाषा है जिसमें !!{function} नहीं हैं। किसी
!!a{structure}~$\Struct M$, !!a{constant}~$c$ और $a \in \Domain M$ के लिए
$\Struct M[a/c]$ को ऐसी !!{structure} परिभाषित करें जो~$\Struct M$ के ठीक
समान है, सिवाय इसके कि $\Assign{c}{M[a/c]} = a$।
$\Struct M \VDash !A$ को !!{sentence}s~$!A$ के लिए इस प्रकार परिभाषित करें:
\begin{enumerate}
\tagitem{prvFalse}{%
  \indcase{!A}{\lfalse}{$\Struct{M} \VDash \indfrm$ नहीं।}}{}

\tagitem{prvTrue}{%
  \indcase{!A}{\ltrue}{$\Struct{M} \VDash \indfrm$।}}{}

\item \indcase{!A}{\Atom{R}{d_1, \dots, d_n}}{$\Struct M \VDash \indfrm$
  तभी और केवल तभी जब $\langle \Assign{d_1}{M}, \dots, \Assign{d_n}{M} \rangle \in
  \Assign{R}{M}$।}

\item \indcase{!A}{\eq[d_1][d_2]}{$\Struct M \VDash \indfrm$ तभी और केवल
  तभी जब $\Assign{d_1}{M} = \Assign{d_2}{M}$।}

\tagitem{prvNot}{%
  \indcase{!A}{\lnot !B}{$\Struct{M} \VDash \indfrm$ तभी और केवल तभी जब
    $\Struct{M} \VDash {!B}$ नहीं।}}{}

\tagitem{prvAnd}{%
  \indcase{!A}{(!B \land !C)}{$\Struct{M} \VDash
    \indfrm$ तभी और केवल तभी जब $\Struct{M} \VDash!B$ और
    $\Struct{M} \VDash !C$।}}{}

\tagitem{prvOr}{%
  \indcase{!A}{(!B \lor !C)}{$\Struct{M} \VDash \indfrm$ तभी और केवल तभी
    जब $\Struct{M} \VDash !B$ अथवा $\Struct{M} \VDash !C$ (अथवा दोनों)।}}{}

\tagitem{prvIf}{%
  \indcase{!A}{(!B \lif !C)}{$\Struct{M} \VDash \indfrm$ तभी और केवल तभी
    जब $\Struct {M} \VDash !B$ नहीं, अथवा $\Struct M \VDash !C$ (अथवा
    दोनों)।}}{}

\tagitem{prvIff}{%
  \indcase{!A}{(!B \liff !C)}{$\Struct{M} \VDash {\indfrm}$ तभी और केवल
    तभी जब या तो $\Struct{M} \VDash {!B}$ और $\Struct{M} \VDash {!C}$,
    दोनों; अथवा न तो $\Struct{M} \VDash {!B}$ और न ही
    $\Struct{M} \VDash {!C}$।}}{}

\tagitem{prvAll}{%
  \indcase{!A}{\lforall[x][!B]}{$\Struct{M} \VDash {\indfrm}$ तभी और केवल
    तभी जब सभी $a \in \Domain{M}$ के लिए
    $\Struct{M[a/c]} \VDash \Subst{!B}{c}{x}$, बशर्ते~$c$,~$!B$ में न
    आता हो।}}{}

\tagitem{prvEx}{%
  \indcase{!A}{\lexists[x][!B]}{$\Struct{M} \VDash
  {\indfrm}$ तभी और केवल तभी जब कोई $a \in \Domain M$ ऐसा हो कि
  $\Struct{M[a/c]} \VDash \Subst{!B}{c}{x}$, बशर्ते~$c$,~$!B$ में न
  आता हो।}}{}
\end{enumerate}
मान लें कि~$x_1$, \dots, $x_n$,~$!A$ के सभी मुक्त !!{variable} हैं;
$c_1$, \dots, $c_n$, ऐसे अचर-चिह्न हैं जो~$!A$ में नहीं आते;
$a_1$, \dots, $a_n \in \Domain M$; और $s(x_i) = a_i$।

यह सिद्ध कीजिए कि $\Sat{M}{!A}[s]$ और निम्न कथन के सत्य-मूल्य समान हैं:

$\Struct M[a_1/c_1,\dots,a_n/c_n]
\VDash \Subst{\Subst{!A}{c_1}{x_1}\dots}{c_n}{x_n}$।

(यह प्रश्न दिखाता है कि प्रथम-क्रम तर्क का ऐसा अर्थविज्ञान देना संभव है जो
चर-नियतनों के बिना काम चला सके।)
\end{prob}

\begin{prob}
मान लें कि~$f$ ऐसा फलन-चिह्न है जो~$!A(x,y)$ में नहीं आता। दिखाइए कि कोई
!!a{structure}~$\Struct{M}$ ऐसी है कि
$\Sat{M}{\lforall[x][\lexists[y][!A(x,y)]]}$ तभी और केवल तभी जब कोई
~$\Struct
M'$ ऐसी है कि $\Sat{M'}{\lforall[x][!A(x,f(x))]}$।

(यह प्रश्न उस परिणाम का एक विशेष प्रकरण है जिसे स्कोलेम का प्रमेय कहते हैं;
$\lforall[x][!A(x,f(x))]$ को
$\lforall[x][\lexists[y][!A(x,y)]]$ का \emph{स्कोलेम सामान्य रूप} कहते हैं।)
\end{prob}

\end{document}
